\section{Methodology}\label{sec:methodology}

\subsection{Dataset}\label{sec:method-dataset}

The dataset has seven non-trivial open-source applications whose maintainers ship their own Kubernetes manifests (Table~\ref{tab:dataset}). A project is eligible if all of the following hold. (i) Its container images are public and pullable without credentials. (ii) It runs on a local \texttt{kind} cluster \citep{tool_kind} without cloud-specific dependencies. (iii) Its manifests may be redistributed; otherwise, we store only the commit. (iv) It is a real application, not a single-container toy. It must have at least two cooperating components, persistent state, or real functionality such as a workflow engine or a storefront. The seven cover four classes and both manifest intents, so the production stratum of RQ1 is not empty. We selected a candidate only if its human manifest could be deployed and measured on the evaluation cluster. So every source of configurations could be scored on all six dimensions. This rule makes the human runtime outcomes partly true by construction. We therefore report them as a reference, not as a test of human against LLM deployability (Section~\ref{sec:res-rq1}).

\begin{table*}[t]
\centering\small
\caption{Evaluated projects. Intent comes from the repository's own wording about its manifests (evidence recorded per project). Size and workloads refer to the shipped human manifest. Stars are those of the repository that holds the manifest, at the fetch date (October 2026).}\label{tab:dataset}
\begin{tabular}{llllrrrr}
\toprule
Project & Class & Intent & Commit & Components & Lines & Objects & Stars \\ \midrule
n8n (n8n-hosting) & stateful & demo & 7a22e91 & 2 & 250 & 9 & 1,750 \\
Grafana & stateful & production & 3920186 & 1 & 79 & 3 & 77,345 \\
Example Voting App & multi-service & demo & 63e9150 & 5 & 182 & 9 & 5,832 \\
Emojivoto & multi-service & demo & 9799218 & 4 & 210 & 11 & 218 \\
Online Boutique & microservice system & demo & 38e7348 & 12 & 912 & 35 & 21,061 \\
Sock Shop & microservice system & demo & 9dff06f & 13 & 844 & 29 & 3,781 \\
Kubernetes Dashboard & cluster add-on & production & 42deb6b & 2 & 242 & 14 & 15,407 \\
\bottomrule
\end{tabular}
\end{table*}

For each project, we build a \texttt{docs\_bundle/} from the project's own README, Dockerfile(s) and Docker Compose file, copied verbatim at the pinned commit. Only Kubernetes-specific content is removed. This covers any embedded manifest fragment, any ``run on Kubernetes'' deployment section, and any single sentence that states a fact otherwise recorded only in the human YAML (e.g., a replica count or a NodePort value). Every redaction is logged with its file, line range and reason. No sentence in \texttt{docs\_bundle/} is written by the research team or by any LLM. An internal \texttt{fact\_sheet.md}, never shown to a model, records the application's purpose and each component's image and required environment variables. It also justifies the per-project requirement checks (D4) and the functional smoke test (D5). Both were written during dataset construction, and every later change to them is one of the logged corrections of Section~\ref{sec:threats}. Before generation, a second author reviews each project's bundle and fact sheet against a fixed checklist.

\subsection{Fairness of the human comparison}\label{sec:method-fairness}

The human developers and the LLMs do not work under the same conditions. The human's goal was often a demo rather than production hardening. The human iterated with real feedback, while each LLM makes a single attempt. The human also had the full repository, while each LLM sees only \texttt{docs\_bundle/}. The design reduces this asymmetry in four ways. First, we treat the human manifest as a real-world reference point. Our claims are comparisons with what the project ships, not with human engineers in general. Second, the human comparison (RQ1) uses the \emph{neutral} prompt. Like the human's situation, this prompt gives no explicit push toward hardening. We analyse the \emph{production} prompt separately (RQ5). Third, we label every human manifest as production- or demonstration-oriented, based on the repository's own wording, and we report RQ1 overall and per label. Fourth, the research team writes no specification. An earlier design used a team-written specification that an LLM had helped to compile. This could have favoured that model family. The final design removes it, so every model receives only the maintainers' own words. The LLM-versus-LLM comparisons (RQ3, RQ5) use identical inputs for every model.

\subsection{Generation}\label{sec:method-generation}

Every model receives the same fixed system prompt, ``You are an experienced DevOps engineer. Respond with Kubernetes YAML manifests only.'' It is followed by a user prompt that joins the files in \texttt{docs\_bundle/} (each preceded by its filename, in a fixed order) and ends in one of two instruction sentences. The \emph{neutral} variant asks to ``write Kubernetes manifests to deploy this application''. The \emph{production} variant asks to ``write production-ready Kubernetes manifests.'' The variants differ in nothing else. We draw three samples per (model, project, prompt variant), each from a fresh context with no memory of any other generation. For each generation, we store the request parameters, the raw response, the extracted YAML and the extraction status. This YAML is every fenced YAML block, joined in order, or the whole response if there is no block. A reply is scored as returned. A reply that opens a YAML block but never closes it falls back to the whole response. So it fails D1, as it would for a user who pasted it.

Table~\ref{tab:models} lists the five models and how each was accessed. For budget reasons, only Gemini and Qwen were accessed through an API (free tiers). Only for them could we set the temperature (0.2) and the output limit (8,192 tokens). The other routes do not expose sampling parameters. So each model is evaluated together with its access route, and RQ3 compares these model-and-access configurations. Claude's outputs came from isolated Claude Code subagents, not a bare API call. The transcripts of all 45 recorded generation sessions show that each subagent read only its own prompt file and wrote only its output manifest. None accessed the human manifests, other models' outputs or the web. These sessions are the 42 retained samples and three pilot samples from an earlier protocol, which were discarded and regenerated. Section~\ref{sec:threats} discusses the consequences of these differences.

\begin{table*}[t]
\centering\small
\caption{Models and access routes. Temperature and output limit could be set only for the API-generated models. For GPT, one author (T.~Hussain) sent the prompts and copied the replies under a written protocol (Section~\ref{sec:threats}).}\label{tab:models}
\resizebox{\textwidth}{!}{%
\begin{tabular}{llllllll}
\toprule
Model & Identifier as recorded & Tier and reasoning & Access route & Generated (2026) & Temp. & Max tokens & Operator \\ \midrule
Claude & claude-sonnet-5-5 & mid tier; harness default & Claude Code subagent, fresh context & 9 Nov & -- & -- & authors \\
DeepSeek & DeepSeek web chat & chat model; reasoning off & chat.deepseek.com, new chat per prompt & 7--9 Nov & -- & -- & authors \\
Gemini & gemini-3.6-flash & flash tier; hidden reasoning & Gemini API (free tier) & 4--11 Nov & 0.2 & 8,192 & authors \\
GPT & ChatGPT default model & default; medium reasoning & chatgpt.com, temporary chat per prompt & 8--13 Nov & -- & -- & T.~Hussain \\
Qwen & qwen3.8-27b & 27B open weights; provider default & Groq API (free tier) & 4--11 Nov & 0.2 & 8,192 & authors \\
\bottomrule
\end{tabular}}
\end{table*}

\subsection{K8sGenBench: a six-dimension benchmark}\label{sec:method-benchmark}

The same automated pipeline evaluates every configuration. It uses deterministic tools (no LLM-as-judge). For D1--D3, it uses at least two tools per dimension.

\textbf{D1: Correctness.} We parse the manifest with PyYAML, rejecting duplicate mapping keys as YAML~1.2 and \texttt{kubectl} do. We then validate it against pinned Kubernetes schemas with \texttt{kubeconform} \citep{tool_kubeconform}. Finally, we submit it object by object to the API server with \texttt{kubectl apply --dry-run=server}. So one object that the cluster cannot resolve does not hide the validity of the others. The metrics are the parse rate, the schema-pass rate and the share of objects the API server accepts.

\textbf{D2: Security.} We scan configurations with Kubescape \citep{tool_kubescape} against the NSA-CISA \citep{nsacisa2022k8shardening} and MITRE frameworks. We report their summary compliance score as \emph{Kubescape compliance}. This score (0--100) is the average, over the controls, of the share of resources that pass each control. We also use Trivy's configuration scanner \citep{tool_trivy} and dry-run the configurations into namespaces that enforce the \texttt{baseline} and \texttt{restricted} Pod Security Standards \citep{k8s_pod_security_standards}. A secret check counts literal values under credential-named keys and known credential token formats, ignoring placeholder values. We also classify findings into the eleven misconfiguration categories of \citet{rahman2023k8smisconfig}, using the rules of their SLI-KUBE tool on every pod and container specification.

\textbf{D3: Static reliability.} We scan configurations with kube-score \citep{tool_kubescore}, Polaris \citep{tool_polaris} and KubeLinter \citep{tool_kubelinter}. This includes the six KubeLinter smells used by \citet{zhang2024genaismells}, reported as density per thousand lines. We also apply a ten-item Reliability Controls Checklist to every workload: at least two replicas for Deployments, a readiness probe, a liveness probe, a CPU request, a memory request, a memory limit, an explicit rolling-update strategy, a PodDisruptionBudget, pod anti-affinity or a topology-spread constraint, and a pinned image tag. Items that do not apply to a workload kind are skipped. Checklist coverage is the share of applicable items met.

\textbf{D4: Requirement satisfaction.} Each project's \texttt{requirements.py} checks facts that can be traced to a file in \texttt{docs\_bundle/}. Examples are that every documented component is deployed or that a documented required environment variable is set. Some facts are true of the project but cannot be recovered from the bundle (for example, an image registry path that appears only in unbundled build scripts). We report these as \emph{oracle} checks and exclude them from the score. They are a diagnostic of what a model knew beyond its input. The metric is the share of scored requirements met. Table~\ref{tab:instruments} (Appendix~\ref{app:fixes}) lists the number of checks and the smoke-test steps of each project.

\textbf{D5: Deployment.} We deploy each configuration into a fresh namespace on a local \texttt{kind} cluster and run a project-specific smoke test: a sequence of HTTP requests from a client pod inside the cluster to the Service in front of the documented entry component. Where the documentation states a data-path effect, the test also runs a command inside a data-store container. Each step expects a documented status code and response content. A configuration counts as \emph{fully Ready} only when every workload has held its desired number of Ready replicas for 20~seconds without interruption. The deadline is 300 seconds. Client pods comply with the \texttt{restricted} Pod Security Standard. They are placed where the configuration's own NetworkPolicies admit a client. So a configuration that deliberately restricts ingress is not mistaken for one that fails to serve. When a configuration does not become Ready, we capture the cluster's record at the deadline. It holds, per container, the restart count, the last exit code and reason and the last log lines, and, per object, the warning events.

\textbf{Root-cause labels.} Each configuration that is not both fully Ready and smoke-test-passing receives one root-cause label from a fixed set of mechanisms (Table~\ref{tab:root-causes}). A label is \emph{observed} when the captured record shows the mechanism directly (for example, an error message that names it). It is \emph{inferred} when we reconstruct it from the manifest together with the observed status. It is \emph{supported} when, in addition, configurations that differ from it only in the suspected setting succeed. One rater assigned the labels from the stored evidence. To check them, a separate instance of an LLM (Claude Sonnet) relabelled every case blind from the same evidence. It ran without access to the first-pass labels or the repository. The evidence was the deployment record and the configuration, with model names removed and cases in random order. The second rater agreed with the first pass on 85.2\% of categories (Cohen's $\kappa=0.84$; \citealp{cohen1960kappa}) and 87.7\% of themes ($\kappa=0.85$). Eleven of the 12 disagreements concerned Sock Shop configurations in which several failures occur together. After we re-examined the evidence, three first-pass labels were changed and nine were kept. The table reports the adjudicated labels. Because the second rater is an LLM, this is a consistency check, not validation by an independent human. The replication package contains the label, mechanism and evidence for every configuration, the blind pack and both passes.

\textbf{D6: Recovery.} We inject two faults into each configuration that serves, three times each, while an in-cluster prober calls the entry Service at 5~Hz. The \emph{pod-delete} fault force-removes a Ready pod of the entry workload. The \emph{container-crash} fault stops its container at the node runtime. The metrics are recovery success, recovery time (from the fault until the replacement pod is Ready), and downtime (failed probes divided by the probing rate). A workload with one replica must be unavailable until its pod is replaced, but one with two or more need not be. So we report D6 by replica count.

We compute no composite score across dimensions.

\subsection{Evaluation pipeline and environment}\label{sec:method-pipeline}

Each configuration passes through the pipeline in a fixed order (Algorithm~\ref{alg:evaluate}): parsing, schema validation, the server-side dry-run, the D2/D3 scanners, the D4 checks, deployment and smoke test (D5), and fault injection (D6). Failures are recorded as data. A parse failure marks every later stage as skipped. A schema failure does not block the static scanners or D4, which read the manifest text. D6 runs when the configuration is fully Ready. Otherwise, it runs when its smoke test passes, because an application can serve correctly while an auxiliary workload (for example, a load generator) is not Ready. We also report D6 for fully Ready configurations only.

Every configuration is deployed into its own namespace, and we recreate the cluster between projects. Every image comes from a local registry mirror filled before deployment, so network latency does not confound time to Ready. Two deployment-time adjustments supply facts that the bundles do not contain. Both are applied identically to every source and recorded per configuration, and no human manifest ever needs them. First, in some projects the exact image path appears only in unbundled build scripts. There, a reference that cannot be pulled as written but whose repository name identifies a documented component is served that component's canonical image. A reference that identifies no documented component fails. Second, an image may enforce a start-up prerequisite that the bundle never states. We supply it when the configuration does not provide it. These prerequisites are two environment variables in Online Boutique (one its front end requires, one that disables a cloud profiler in two services), and an empty Secret with the fixed name that the Kubernetes Dashboard expects. Permissions are never added. Section~\ref{sec:threats} bounds the effect of both adjustments.

The cluster has one control-plane node and three worker nodes. So configurations that use pod anti-affinity or a topology spread across three replicas can schedule. One setting differs from \texttt{kind}'s defaults. The container runtime limits each container to 1,048,576 open files (Docker's default). We set this because some older images size internal tables from that limit and could otherwise exhaust the host's memory. Image tags are resolved when the mirror is filled, at evaluation time. A \texttt{LoadBalancer} Service is treated as valid but pending. An Ingress is created but not exercised. PersistentVolumeClaims use \texttt{kind}'s local-path StorageClass. So a configuration that requests a StorageClass absent from this environment fails, and the reason is recorded.

\begin{algorithm}[t]
\caption{Evaluating one configuration $c$ of project $p$.}\label{alg:evaluate}
\small
\begin{algorithmic}[1]
\State $r \gets$ new result record for $c$
\State $(docs, text, err) \gets \Call{StrictParse}{c}$ \Comment{duplicate keys rejected}
\If{$err \neq \varnothing$}
  \State mark D1--D6 of $r$ as skipped; \Return $r$
\EndIf
\State $r.\mathrm{D1} \gets \Call{Kubeconform}{text}$, \Call{ServerDryRun}{$docs$} \Comment{object by object}
\State $r.\mathrm{D2}, r.\mathrm{D3} \gets \Call{Scanners}{docs, text}$, \Call{PodSecurityDryRun}{$docs$}
\State $r.\mathrm{D4} \gets \Call{Requirements}{p, docs}$
\State $ns \gets$ fresh namespace; serve images from the local mirror; apply $docs$ to $ns$
\State $t_0 \gets now$; $since \gets \bot$
\While{$now - t_0 < 300$\,s}
  \If{every workload has its desired Ready replicas}
    \State $since \gets since$ \textbf{or} $now$
    \If{$now - since \geq 20$\,s}
      \State \textbf{break}
    \EndIf
  \Else
    \State $since \gets \bot$
  \EndIf
  \State wait 3\,s
\EndWhile
\State $r.\mathrm{ready} \gets (since \neq \bot)$; \textbf{if not} $r.\mathrm{ready}$ \textbf{then} capture diagnostics
\State $r.\mathrm{smoke} \gets \Call{SmokeTest}{p, ns}$ \Comment{client placed where NetworkPolicies admit it}
\If{$r.\mathrm{ready}$ \textbf{or} $r.\mathrm{smoke}$ passed}
  \State $r.\mathrm{D6} \gets$ 3 pod deletes and 3 container crashes under a 5\,Hz prober
\EndIf
\State delete $ns$ \Comment{always, even after an error}
\State \Return $r$
\end{algorithmic}
\end{algorithm}

\subsection{Statistical analysis}\label{sec:method-analysis}

We report medians for scores and percentages for pass/fail outcomes. Because the three samples of a (model, project, prompt) cell share their input, and the projects differ widely, our tests treat the project or the (model, project) cell as the unit wherever possible. We use rank-based tests: Kruskal--Wallis and Fisher's exact test to compare sources (RQ1), a Friedman test blocked on project to compare models (RQ3), a logistic regression to relate static scores to deployment success (RQ4), and a Wilcoxon signed-rank test to compare the two prompts within each (model, project) cell (RQ5). $P$-values are Holm-adjusted for multiple comparisons. Appendix~\ref{app:stats} describes every test, the effect sizes and the robustness checks, and reports their results.
